\subsection{流形的例子}

5.2.1. 设 X 是集合；X 上存在唯一的流形结构，使其底层拓扑空间是离散的；该结构是维数为 0 的纯流形结构。

5.2.2. 设 E 是巴拿赫空间。三元组 \(c = (\mathrm{E}, \mathrm{Id}_{\mathrm{E}}, \mathrm{E})\) 是 E 的坐标图，且 \(\mathcal{A} = \{c\}\) 是 E 的图册，因而在 E 上定义 E 型纯流形结构；底层拓扑是 E 上给定的拓扑。以下谈及 E 上的流形结构时，始终指上述结构。

特别地，这适用于每个赋予与其向量结构相容的唯一豪斯多夫拓扑的 K 上有限维向量空间（Esp. Vect. Top.，第 I 章，第 2 节，第 3 号）。

5.2.3. 设 X 是流形，U 是 X 的开子集。U 上存在一种流形结构，其坐标图是定义域包含于 U 的流形 X 的坐标图。称此结构由 X 的结构诱导（参见第 5.3.4 号）；赋予此结构的 U 称为 X 的开子流形。

特别地，巴拿赫空间 E 的每个开子集都赋有 E 型纯流形的规范结构。设 X 是流形；三元组 (U，\(\varphi\)，E) 是 X 的坐标图的充要条件是 U 在 X 中开，且 \(\varphi\) 是从 X 的开子流形 U 到 E 的开子流形的同构。

5.2.4. 设 X 是集合，\((X_{i})_{i\in I}\) 是 X 的一个覆盖。假设每个 \(X_{i}\) 上给定一种流形结构，满足下列条件：

对 I 中任意 \(i\) 和 \(j\)，集合 \(\mathbf{X}_{i} \cap \mathbf{X}_{j}\) 在 \(\mathbf{X}_{i}\) 和 \(\mathbf{X}_{j}\) 中都是开集，并且由 \(\mathbf{X}_{i} \cap \mathbf{X}_{j}\) 与 X＿j 的结构在 \(\mathbf{X}_{i}\)  \(\mathbf{X}_{j}\) 上诱导的流形结构相同。

于是 X 上存在唯一的流形结构，使得对 I 中每个 i，\(X_{i}\) 都是 X 的开子流形。称流形 X 由各流形 \(X_{i}\) 粘合而成。

5.2.5. 设 X 是流形。X 中满足 \(X_{n}\) 的点 x 的集合 \(\dim_{x}X=n\ (n\ \text{为整数}\geqslant0)\) 是 X 的开子流形，且是纯 n 维的。

5.2.6. 设 E 是巴拿赫空间。可以按如下方式在具有拓扑补的 E 的向量子空间集合 \(\mathbf{G}(\mathbf{E})\) 上赋予解析流形结构：对每一对 \((\mathrm{F}_0, \mathrm{G}_0) \in \mathbf{G}(\mathbf{E}) \times \mathbf{G}(\mathbf{E})\)，若 \(\mathrm{E} = \mathrm{F}_0 \oplus \mathrm{G}_0\)，用 \(\mathrm{U}_{\mathrm{G}_0}\) 表示具有 \(\mathrm{F} \in \mathbf{G}(\mathbf{E})\) 作为补空间的 \(\mathrm{G}_0\) 的集合，并通过将每个 \(\varphi_{\mathrm{F}_0, \mathrm{G}_0}\) 与从 \(\mathrm{U}_{\mathrm{G}_0}\) 到 \(\mathcal{L}(\mathrm{F}_0; \mathrm{G}_0)\) 的映射关联起来定义从 \(\mathrm{F} \in \mathrm{U}_{\mathrm{G}_0}\) 到巴拿赫空间 \(\mathrm{F}_0\) 的双射 \(\mathrm{G}_0\)，该映射的图像是 \(\mathrm{E} = \mathrm{F}_0 \times \mathrm{G}_0\) 的子空间 F。坐标图 \((\mathrm{U}_{\mathrm{G}_0}, \varphi_{\mathrm{F}_0, \mathrm{G}_0}, \mathcal{L}(\mathrm{F}_0; \mathrm{G}_0))\) 构成 \(\mathbf{G}(\mathbf{E})\) 的一个图册。赋予由此图册定义的流形结构后，\(\mathbf{G}(\mathbf{E})\) 称为 E 的格拉斯曼流形。

拓扑空间 G(E) 可度量化。若 K 局部紧且 E 有限维，则 G(E) 紧。

对每个整数 \(n \geqslant 0\)，由维数（分别地，余维）为 \(\mathbf{G}_{n}(\mathbf{E})\) 的 \(\mathbf{G}^{n}(\mathbf{E})\) 构成的集合 \(\mathbf{F} \in \mathbf{G}(\mathbf{E})\)（分别地，G＾\(n\)(E)）在 \(\mathbf{G}(\mathbf{E})\) 中既开又闭，并且是 \(\mathbf{G}(\mathbf{E})\) 的纯开子流形。若 \(\mathbf{K} = \mathbf{R}\) 或 \(\mathbf{C}\)，则对每个 n，\(\mathbf{G}_{n}(\mathbf{E})\)\(n\) 都是连通的。

将每个 \(F \in \mathbf{G}(\mathbf{E})\) 关联到其在 E 的对偶 \(E'\) 中的正交补的映射，是从 \(\mathbf{G}(\mathbf{E})\) 到 \(\mathbf{G}(\mathbf{E}')\) 的态射，并对每个整数 n 诱导出从 \(\mathbf{G}^{n}(\mathbf{E})\) 到 \(\mathbf{G}_{n}(\mathbf{E}')\) 的同构。

若 K = R 或 C，或者 E 有限维，则 \(\mathbf{G}_{1}(\mathbf{E})\) 正是与 E 关联的射影空间，因而赋有解析流形结构。

